% Part: incompleteness
% Chapter: introduction
% Section: historical-background

\documentclass[../../../include/open-logic-section]{subfiles}

\begin{document}

\olfileid{inc}{int}{bgr}

\olsection{ઐતિહાસિક પૃષ્ઠભૂમિ}

આ વિભાગમાં અપૂર્ણતાનાં પ્રમેયોને યોગ્ય સંદર્ભમાં સમજવામાં
મદદરૂપ થતી ઐતિહાસિક ઘટનાઓની ટૂંકી ચર્ચા કરીશું.
ખાસ કરીને, ગાણિતિક તર્કશાસ્ત્રના ઇતિહાસનો ખૂબ જ
સંક્ષિપ્ત આલેખ આપીશું અને પછી ગણિતના પાયા વિશેના
ઇતિહાસની થોડી વાત કરીશું.

\begin{digress}
  ``ગાણિતિક તર્કશાસ્ત્ર'' શબ્દસમૂહના બે અર્થ થઈ શકે.
  ``ગાણિતિક'' શબ્દ વિષયને વર્ણવે છે એમ લઈએ, તો
  ``ગણિતનો તર્ક''નો અર્થ થાય: ગણિતીય દલીલોના સિદ્ધાંતો.
  એ શબ્દ પદ્ધતિને વર્ણવે છે એમ લઈએ, તો ``તર્કનું ગણિત''નો
  અર્થ થાય: દલીલોના સિદ્ધાંતોનો ગાણિતિક અભ્યાસ.
  આગળનું વર્ણન ઘણી વાર એકસાથે આ બંને અર્થમાં ગાણિતિક
  તર્કશાસ્ત્ર વિશે છે.
\end{digress}

તર્કશાસ્ત્રનો અભ્યાસ મૂળે એરિસ્ટોટલથી શરૂ થયો; તે
લગભગ 384--322 \textsc{bce} દરમિયાન જીવ્યા.
તેમનાં \emph{Categories}, \emph{Prior
  analytics} અને \emph{Posterior analytics}માં
વૈજ્ઞાનિક દલીલોના સિદ્ધાંતોનો વ્યવસ્થિત અભ્યાસ છે;
તેમાં નિગમનત્રયનો પણ વિગતવાર અને વ્યવસ્થિત અભ્યાસ છે.

મધ્યયુગના પંડિતીય દર્શનમાં એરિસ્ટોટલનું તર્કશાસ્ત્ર
પ્રભુત્વ ધરાવતું હતું. અઢારમી સદી જેટલા મોડા સમયમાં પણ
કાન્ટ માનતા હતા કે એરિસ્ટોટલનું તર્કશાસ્ત્ર સંપૂર્ણ છે
અને તેમાં સુધારાની જરૂર નથી. પરંતુ નિગમનત્રયનો સિદ્ધાંત
ગણિતીય દલીલના માત્ર ઉપરછલ્લા પાસાંઓનું જ નિરૂપણ કરી
શકે છે. તેની એક સદી પહેલાં, ન્યૂટનના સમકાલીન લાઇબનિટ્ઝે
તાર્કિક દલીલ માટે સંપૂર્ણ ``કલન''ની કલ્પના કરી હતી અને
એવું કલન રચવા તરફ કેટલાક પ્રાથમિક પગલાં પણ લીધાં;
તેમણે મૂળે વિધાનાત્મક તર્કશાસ્ત્રનું એક સ્વરૂપ વર્ણવ્યું.

ઓગણીસમી સદી તર્કશાસ્ત્ર માટે નિર્ણાયક બની. 1854માં
જ્યોર્જ બૂલે \emph{The Laws of Thought} લખ્યું;
તેમાં વિધાનાત્મક તર્કશાસ્ત્રનો વિસ્તૃત બીજગણિતીય અભ્યાસ
છે, જે આજની રજૂઆતોથી બહુ દૂર નથી. 1879માં ગોટલોબ
ફ્રેગેએ \emph{Begriffsschrift} (સંકેતલેખન) પ્રકાશિત
કર્યું. તે વિધાનાત્મક તર્કશાસ્ત્રમાં પરિમાણકો અને સંબંધો
ઉમેરીને પ્રથમ-ક્રમનું તર્કશાસ્ત્ર પણ સમાવે છે. વાસ્તવમાં,
ફ્રેગેની તાર્કિક પદ્ધતિઓમાં ઉચ્ચ-ક્રમનું તર્કશાસ્ત્ર અને
તેથી વધુ પણ હતું. પોતાના \emph{Basic Laws of
  Arithmetic}માં તેમણે બતાવવાનો પ્રયત્ન કર્યો કે
અંકગણિતનું સમગ્ર જ્ઞાન કેવળ તાર્કિક પૂર્વધારણાઓમાંથી
તેમના \emph{Begriffsschrift}માં વ્યૂત્પન્ન થઈ શકે.
દુર્ભાગ્યે, 1902માં રસેલે બતાવ્યું તેમ આ પૂર્વધારણાઓ
અસુસંગત નીવડી. તે અસુસંગત સ્વયંસિદ્ધ વાક્ય બાજુએ
રાખીએ તો ફ્રેગેએ લગભગ એકલા હાથે આધુનિક તર્કશાસ્ત્ર
શોધ્યું—આ નોંધપાત્ર સિદ્ધિ છે. બૂલ પછી બીજગણિતીય
અભિગમ ધરાવતા પીયર્સ અને શ્રોડર સહિતના વિચારકોએ પણ
પરિમાણકયુક્ત તર્કશાસ્ત્ર સ્વતંત્ર રીતે વિકસાવ્યું.

હવે ગણિતના પાયા અંગેના વિકાસ તરફ વળીએ. ગણિતમાં
તર્કશાસ્ત્રની મહત્ત્વની ભૂમિકા હોવાથી ઉપર વર્ણવેલા
વિકાસો સાથે તેનો ઘણો સંબંધ છે. દાખલા તરીકે, સમગ્ર
ગણિત માત્ર પોતાની તાર્કિક ગોઠવણી પર આધારિત છે એવું
બતાવવાના સ્પષ્ટ હેતુથી ફ્રેગેએ તેમનું તર્કશાસ્ત્ર
વિકસાવ્યું. ખાસ કરીને, કાન્ટના મત મુજબનાં અનુભવથી
સ્વતંત્ર \emph{સંશ્લેષક} સત્યોને બદલે ગણિત અનુભવથી
સ્વતંત્ર \emph{વિશ્લેષક} સત્યોનું બનેલું છે એમ તેઓ બતાવવા ઇચ્છતા હતા.

ઘણા લોકો શુદ્ધ ગણિતનો જન્મ ગ્રીક લોકો પાસે થયો એમ
માને છે. લગભગ 300 B.C.માં લખાયેલું યુક્લિડનું
\emph{Elements} ચોકસાઈ અને કડક દલીલ પર ભાર મૂકતું
ગ્રીક ગણિતનું પરિપક્વ ઉદાહરણ છે. તેના વ્યાખ્યાઓ અને
સાબિતીઓ આજના માધ્યમિક શાળાના ભૂમિતિનાં પાઠ્યપુસ્તકોમાં
લગભગ મૂળ સ્વરૂપે જ છે (જ્યાં સુધી ભૂમિતિ હજી શીખવાય છે).
ગણિતમાં જ નહીં, દર્શનની શાખાઓમાં પણ આ ગાણિતિક દલીલની
રીતને કડક દલીલનું આદર્શ માનવામાં આવ્યું છે. સ્પિનોઝાએ તો
નૈતિક અને ધાર્મિક દલીલો પણ યુક્લિડની શૈલીમાં રજૂ કરી હતી—
જે જોવામાં વિચિત્ર લાગે છે!

સત્તરમી સદીમાં ન્યૂટન અને લાઇબનિટ્ઝે કલન શોધ્યું.
(સદીઓ સુધી પ્રાથમિકતાનો ઉગ્ર વિવાદ રહ્યો, પરંતુ આજે
મોટાભાગના વિદ્વાનો માને છે કે બંને વિકાસ મુખ્યત્વે
સ્વતંત્ર હતા.) કલનમાં, દાખલા તરીકે, અનંત સૂક્ષ્મ
રાશિઓના અનંત સરવાળા વિશે દલીલો આવે છે. આ લક્ષણોને
કારણે બિશપ બર્કલીએ તેની ટીકા કરી: તેમના મતે ઈશ્વરમાં
માનવું તેમના સમયના ગણિત કરતાં ઓછું તર્કસંગત નહોતું.
અઢારમી સદીમાં કલનની પદ્ધતિઓનો વ્યાપક ઉપયોગ થયો;
લિયોનાર્ડ ઓયલરે અનંત સરવાળાઓવાળી ગણતરીઓથી નોંધપાત્ર
પરિણામો મેળવ્યાં.

ઓગણીસમી સદીમાં ગણિતશાસ્ત્રીઓએ કલનને વધુ મજબૂત પાયો
આપીને બર્કલીની ટીકાનો જવાબ આપવાનો પ્રયત્ન કર્યો.
કોશી, વાયરસ્ટ્રાસ, બોલ્ઝાનો વગેરેના પ્રયાસોથી
``એપ્સિલોન અને ડેલ્ટા'' વડે લક્ષ, સાતત્ય,
અવકલન અને સંકલનની આજની વ્યાખ્યાઓ મળી; એટલે કે
તેમાં અનંતસૂક્ષ્મ રાશિઓનો કોઈ ઉલ્લેખ નથી. સદીના
પછીના ભાગમાં ગણિતશાસ્ત્રીઓ આગળ વધીને વાસ્તવિક
સંખ્યાઓ સહિત સમગ્ર કલનને પ્રાકૃતિક સંખ્યાઓની મદદથી
સમજાવવા પ્રયત્નશીલ બન્યા. (ક્રોનેકર: ``પૂર્ણ સંખ્યાઓ
ઈશ્વરે બનાવી; બાકી બધું માણસનું કામ છે.'') 1872માં
ડેડકિન્ડે ``Continuity and the irrational numbers''
લખ્યું. તેમાં તેમણે સંમેય સંખ્યાઓના ગણ તરીકે વાસ્તવિક
સંખ્યાઓ કેવી રીતે ``રચવી'' તે બતાવ્યું (અને સંમેય
સંખ્યાઓને, આપણે જાણીએ છીએ તેમ, પ્રાકૃતિક સંખ્યાઓની
જોડીઓ તરીકે જોઈ શકાય). 1888માં તેમણે ``Was sind und
was sollen die Zahlen'' (લગભગ, ``પ્રાકૃતિક સંખ્યાઓ
શું છે અને તે કેવી હોવી જોઈએ?'') લખ્યું; તેમાં
પ્રાકૃતિક સંખ્યાઓને કેવળ ``તાર્કિક'' શબ્દોમાં સમજાવવાનો
હેતુ હતો. 1887માં ક્રોનેકરે ``\"Uber den Zahlbegriff''
(``સંખ્યાની સંકલ્પના વિશે'') લખ્યું; તેમાં તેમણે તમામ
ગણિતીય પદાર્થોને પૂર્ણ સંખ્યાઓ વડે નિરૂપવાની વાત કરી.
1889માં જ્યુસેપ્પે પિયાનોને પ્રાકૃતિક સંખ્યાઓ માટે
ઔપચારિક, સાંકેતિક સ્વયંસિદ્ધ વાક્યો આપ્યાં.

ઓગણીસમી સદીના અંતે અનંત અંગે વધુ સાહસિક અભિગમ પણ
દેખાયો. તે પહેલાં, અનંત પદાર્થો અને સંરચનાઓને
(જેમ કે પ્રાકૃતિક સંખ્યાઓનો ગણ) સાવચેતીથી જોવામાં
આવતાં: ``અનંત સંખ્યામાં'' એટલે ``જેટલાં જોઈએ તેટલાં'',
અને ``લક્ષ તરફ જાય'' એટલે ``જેટલા જોઈએ તેટલા નજીક
જાય'' એવો અર્થ લેવાતો. પરંતુ જ્યોર્જ કૅન્ટૉરે બતાવ્યું
કે અનંતને ખરેખર અનંત તરીકે લઈ શકાય. કૅન્ટૉર,
ડેડકિન્ડ અને અન્યના કાર્યથી આજે વ્યાપક સ્વીકૃત
ગણસિદ્ધાંત આધારિત ગણિતની સામાન્ય સમજ વિકસી.

આથી આપણે વીસમી સદીના તર્કશાસ્ત્ર અને પાયાના વિકાસો
સુધી પહોંચીએ છીએ. 1902માં રસેલે ફ્રેગેની તાર્કિક
પદ્ધતિમાં વિરોધાભાસ શોધ્યો. 1904માં ઝર્મેલોએ
કહેવાતા ``પસંદગીના સ્વયંસિદ્ધ વાક્ય'' વડે કૅન્ટૉરનો
સુવ્યવસ્થિત-ક્રમનો સિદ્ધાંત સાબિત કર્યો; આ સ્વયંસિદ્ધ
વાક્યની માન્યતા અંગે ઘણો વિવાદ થયો. 1910થી 1913
વચ્ચે રસેલ અને વ્હાઇટહેડના \emph{Principia Mathematica}ના
ત્રણ ખંડ પ્રસિદ્ધ થયા; તેમાં ગણિતને તર્કના પાયા પર
ઊભું કરવાના ફ્રેગેના કાર્યક્રમને આગળ વધાર્યો. પરંતુ
રસેલ અને વ્હાઇટહેડે બે એવા સિદ્ધાંતો સ્વીકારવા પડ્યા
જેને કેવળ તાર્કિક ગણવાનું મુશ્કેલ લાગતું હતું:
અનંતનું સ્વયંસિદ્ધ વાક્ય અને ``ઘટાડનીયતા''નું
સ્વયંસિદ્ધ વાક્ય. 1900ના દાયકામાં પોઇનકારે ગણિતમાં
``પોતે વ્યાખ્યાયિત થનારને સમાવતી સમષ્ટિ પર આધારિત
વ્યાખ્યાઓ''ના ઉપયોગની ટીકા કરી;
1910ના દાયકામાં બ્રાઉવરે સમગ્ર ગણિતને ``સ્ફુરણાવાદી''
પાયા પર નવેસરથી સ્થાપવાનો પ્રસ્તાવ મૂક્યો; તેમાં
વર્જિત મધ્યનો નિયમ ($!A \lor \lnot !A$) વપરાતો નથી.

ખરેખર વિચિત્ર સમય!{} સમગ્ર ગણિતને તર્કમાં ઘટાડવાનો
કાર્યક્રમ હવે ``તર્કવાદ'' કહેવાય છે અને ઉપર જણાવેલી
મુશ્કેલીઓને કારણે તેને સામાન્ય રીતે નિષ્ફળ માનવામાં
આવે છે. માનસિક સ્ફુરણાત્મક રચનાઓ વડે ગણિત વિકસાવવાનો
કાર્યક્રમ ``સ્ફુરણાવાદ'' કહેવાય છે; રોજિંદા ગણિત પર
તે અતિ કડક મર્યાદાઓ મૂકે છે એવું મનાય છે. સદીના
વળાંકે ડેવિડ હિલ્બર્ટ—અત્યંત પ્રભાવશાળી ગણિતશાસ્ત્રીઓમાંના
એક—કૅન્ટૉર અને ડેડકિન્ડે રજૂ કરેલી નવી અમૂર્ત
પદ્ધતિઓના પ્રબળ સમર્થક હતા: ``કૅન્ટૉરે આપણા માટે
રચેલા સ્વર્ગમાંથી આપણને કોઈ હાંકી કાઢશે નહીં.''
તેમ છતાં તેઓ આ નવી પદ્ધતિઓની પાયા અંગેની ટીકાઓને
પણ ગંભીરતાથી લેતા (વિચિત્ર રીતે, હવે આ પદ્ધતિઓ
``શાસ્ત્રીય'' કહેવાય છે). તેમણે બંને હેતુ સાધવાનો
નીચેનો રસ્તો સૂચવ્યો:
\begin{enumerate}
\item શાસ્ત્રીય પદ્ધતિઓને ઔપચારિક સ્વયંસિદ્ધ વાક્યો અને
  નિયમો વડે નિરૂપવી; ગણિતીય પ્રશ્નોને સ્વયંસિદ્ધ-તંત્રમાં
  !!{formula}s તરીકે નિરૂપવા.
\item સલામત, ``સાન્ત-પદ્ધતિઓ'' વડે આ ઔપચારિક
  નિષ્પત્તિ-તંત્રો સુસંગત છે તે સાબિત કરવું.
\end{enumerate}

પહેલો હેતુ સિદ્ધ કરવા હિલ્બર્ટનું કાર્ય ઘણું આગળ
વધ્યું. 1899માં તેમના પ્રસિદ્ધ ગ્રંથ \emph{Foundations
of geometry}માં તેમણે ભૂમિતિ માટે આવું કર્યું હતું.
પછીનાં વર્ષોમાં તેમણે અને તેમના વિદ્યાર્થીઓ તથા
સહયોગીઓએ ગણિતની બીજી શાખાઓમાં પણ એવું કરવાનો
પ્રયત્ન કર્યો. હિલ્બર્ટે પોતે અંકગણિત અને વિશ્લેષણ
માટે સ્વયંસિદ્ધ-તંત્રો આપ્યાં. ઝર્મેલોએ ગણસિદ્ધાંતનું
સ્વયંસિદ્ધીકરણ આપ્યું; ફ્રેન્કેલ, સ્કોલેમ,
ફોન નોયમાન અને અન્યોએ તેને વિસ્તાર્યું. 1920ના
દાયકાના મધ્ય સુધી ``સમગ્ર'' ગણિતના સ્વયંસિદ્ધીકરણનો
દાવો કરતા બે અભિગમ હતા: રસેલ અને વ્હાઇટહેડનું
\emph{Principia mathematica} અને જે ઝર્મેલો--ફ્રેન્કેલ
ગણસિદ્ધાંત કહેવાયું તે.

1921માં હિલ્બર્ટે આ તંત્રોની સુસંગતતા સાબિત કરવાનો
હેતુ ધરાવતો સંશોધન-કાર્યક્રમ શરૂ કર્યો. તેમાં તેમના
કેટલાક વિદ્યાર્થીઓ, ખાસ કરીને બર્નેય્સ, એકરમાન
અને પછી જેન્ટ્ઝેને મદદ કરી. મુખ્ય વિચાર આ હતો:
ગણિતમાં અસુસંગતતાની !!{derivation} શક્ય છે કે નહીં
તે પ્રશ્નને ચિહ્નોની શક્ય શ્રેણીઓ અંગેના સંયોજનાત્મક
પ્રશ્નમાં ફેરવવો; એટલે કે અંકગણિત, વિશ્લેષણ કે
ગણસિદ્ધાંતના સ્વયંસિદ્ધ-તંત્રમાંથી, દાખલા તરીકે,
$!A \land \lnot !A$ની સાચી !!{derivation} હોવાના
માપદંડને સંતોષતી વાક્યોની શક્ય શ્રેણીઓનો પ્રશ્ન.
આવી ચિહ્ન-શ્રેણી શક્ય નથી તેની સાબિતી પોતે ગાણિતિક
સાબિતી હોવાથી આ સ્વયંસિદ્ધ-તંત્રોમાં ઔપચારિક બનાવી
શકાય. બીજા શબ્દોમાં, દાખલા તરીકે અંકગણિત સુસંગત
છે એવું કહેતો કોઈ વાક્ય~$\OCon$ હશે. વધુમાં,
ખાસ કરીને જો તેની સાબિતી માટે માત્ર ખૂબ મર્યાદિત,
``સાન્ત'' સાધનો જોઈએ, તો તે વાક્ય સંબંધિત તંત્રોમાં
સાબિત થઈ શકવું જોઈએ.

બીજો હેતુ—વિકસાવેલા સ્વયંસિદ્ધ-તંત્રો દરેક ગણિતીય
પ્રશ્ન નક્કી કરે—બે રીતે ચોક્કસ કરી શકાય. એક રીતે
આવું કહી શકાય: ગણિતના સ્વયંસિદ્ધ-તંત્રની ભાષાનાં
દરેક વાક્ય~$!A$ માટે સ્વયંસિદ્ધ વાક્યોમાંથી
કાં તો~$!A$ અથવા~$\lnot !A$ સાબિત થાય.
આ સાચું હોય તો એવું કોઈ વાક્ય નહીં હોય જેને
આ સ્વયંસિદ્ધ વાક્યોના આધારે ન સાબિત કરી શકાય
કે ન નકારી શકાય; તંત્ર કોઈ પ્રશ્ન અનિર્ણીત નહીં
રાખે. આ ગુણધર્મવાળું સ્વયંસિદ્ધ-તંત્ર
\emph{પૂર્ણ} કહેવાય છે. અલબત્ત, કોઈ ચોક્કસ
વાક્ય માટે બેમાંથી કઈ શક્યતા છે તે નક્કી કરવું
હજી મુશ્કેલ હોઈ શકે. પરંતુ સિદ્ધાંતમાં એવું
કરવાની પદ્ધતિ હોવી જોઈએ. વાસ્તવમાં, હિલ્બર્ટે
વિચારેલાં સ્વયંસિદ્ધ અને !!{derivation} તંત્રો માટે
પૂર્ણતા એવી પદ્ધતિનું અસ્તિત્વ ફલિત કરે—જોકે
હિલ્બર્ટ આ વાત જાણતા નહોતા. પ્રશ્નનો બીજો
અર્થ વધુ મજબૂત શરત છે: આપેલું વાક્ય~$!A$
સ્વયંસિદ્ધ વાક્યોમાંથી વ્યૂત્પન્ન થાય છે કે નહીં
તે નક્કી કરતી યાંત્રિક, સંગણકીય પદ્ધતિ હોય.

1931માં ગડેલે બે ``અપૂર્ણતાનાં પ્રમેયો'' સાબિત
કર્યાં, જે બતાવે છે કે આ કાર્યક્રમ સફળ થઈ
શકે નહીં. ગણિતના પૂરતા શક્તિશાળી,
સુસંગત અને સ્વયંસિદ્ધીકરણીય તંત્રો પૂર્ણ
હોઈ શકતાં નથી. વધુ શરતો હેઠળ આવું તંત્ર
પોતાની સુસંગતતા વ્યક્ત કરતું વાક્ય સાબિત
કરી શકતું નથી.
\sourcecorrection{OLINC-001}{સ્થિર અંગ્રેજી મૂળ અહીં
``ગણિત માટેનું કોઈ તંત્ર પૂર્ણ નથી'' અને સુસંગતતાનું
વાક્ય ``ન સાબિત થાય કે ન નકારી શકાય'' એમ અતિ વ્યાપક
રીતે કહે છે. આ જ ગ્રંથના અપૂર્ણતા-પરિણામોના અવલોકનમાં
પ્રથમ પ્રમેય માટે સુસંગત, સ્વયંસિદ્ધીકરણીય અને
સંગણનીય વિધેયો/નિર્ણેય સંબંધોને નિરૂપતું તંત્ર
ધાર્યું છે; બીજા પ્રમેયમાં પોતાની સુસંગતતા સાબિત ન
થવાની વાત વધુ કડક શરતો હેઠળ છે. મૂળનું નકારી ન
શકાય તેવું સામાન્ય દાવું અહીં રાખ્યું નથી.}

આ હિલ્બર્ટના મૂળ કાર્યક્રમ માટે ભારે આઘાત હતો.
તેમ છતાં, ગણિતમાં ઘણી વાર થાય છે તેમ, તેનાથી
સંશોધન માટે નવી અને રસપ્રદ દિશાઓ પણ ખુલી.
ગણિતનું એક જ સર્વસમાવેશક ઔપચારિક તંત્ર ન હોય,
તો વધુ મર્યાદિત તંત્રો વિકસાવી તેમાં શું સાબિત
કરી શકાય તે તપાસવું યોગ્ય છે. આ તંત્રોની
સુસંગતતા સ્થાપવા માટે ઓછા મર્યાદિત સાબિતી-માર્ગો
વિકસાવવાનું અને તેમની સુસંગતતા સાબિત કરવી કેટલી
મુશ્કેલ છે તે માપવાનું પણ યોગ્ય છે. ગડેલે બતાવ્યું
કે (લગભગ) દરેક ઔપચારિક તંત્રમાં તે નક્કી ન
કરી શકે એવા પ્રશ્નો છે; તેથી કોઈ ચોક્કસ તંત્ર
નક્કી ન કરી શકે એવા ``રસપ્રદ'' પ્રશ્નો શોધવા
અને તેમને નક્કી કરવા માટે તંત્ર કેટલું શક્તિશાળી
હોવું જોઈએ તે સમજવું યોગ્ય છે. આજ સુધી
તર્કશાસ્ત્રીઓ સાબિતીઓના સિદ્ધાંત નામની નવી
ગણિતીય શાખામાં આ પ્રશ્નોનો અભ્યાસ કરે છે.

\end{document}
